% !TeX TS-program = xelatex
\documentclass[aspectratio=169, 12pt, xcolor=table]{beamer}
\usefonttheme{professionalfonts}
\usefonttheme{serif}
\usepackage[T1]{fontenc}
\usepackage{fontspec-xetex}

\usepackage{booktabs}
\usepackage{listings}
\usepackage{subcaption}

\setmainfont{Lato}

%\PassOptionsToPackage[more=table]{xcolor}

% Local configuration
\renewcommand{\figurename}{}
\DeclareCaptionFormat{custom}
{%
	\tiny #3
}
\captionsetup{format=custom}

% Title stuff
\title{Statistics}
\subtitle{An Introduction using R}
\date{Week 01}
\author{Vassilis Markos, Mediterranean College}

\usetheme{streamline}
%\usepackage{/theme/beamerthemestreamline}

% Local Commands
\newcommand{\ohref}[1]{\href{#1}{\texttt{#1}}}
\def\sms#1{\smsgo#1\relax}
\def\smsgo#1.#2\relax{#1.{\scriptsize #2}}

% Code listings

\definecolor{codegreen}{rgb}{0,0.6,0}
\definecolor{codegray}{rgb}{0.5,0.5,0.5}
\definecolor{codepurple}{rgb}{0.58,0,0.82}
\definecolor{backcolour}{rgb}{0.95,0.95,0.92}

\lstdefinestyle{mystyle}{
	backgroundcolor=\color{backcolour},   
	commentstyle=\color{codegreen},
	keywordstyle=\color{magenta},
	numberstyle=\tiny\color{codegray},
	stringstyle=\color{codepurple},
	basicstyle=\ttfamily\footnotesize,
	breakatwhitespace=false,         
	breaklines=true,                 
	captionpos=b,                    
	keepspaces=true,                 
	numbers=left,                    
	numbersep=5pt,                  
	showspaces=false,                
	showstringspaces=false,
	showtabs=false,                  
	tabsize=2
}

\lstset{style=mystyle}

\begin{document}

	\begin{frame}
		\titlepage
	\end{frame}

	\begin{frame}{Contents}
		\tableofcontents
	\end{frame}

	\section{Descriptive Statistics}\label{sec:descriptive-statistics}
	
	\sectionframe
	
	\begin{frame}{Stopwatch!}
		Open your mobile phone stopwatch application or just use:
		\begin{center}
			\url{https://www.online-stopwatch.com/}
		\end{center}
		Then:
		\begin{itemize}
			\item Start your stopwatch and try to press the "Stop!" (or "Pause") button at exactly \sms{3.000} seconds.
			\item Repeat the above process 30 times, noting down the actual time shown when you stopped your stopwatch.
		\end{itemize}
		Okay, now, who's the best reflexes in this class?
	\end{frame}
	
	\begin{frame}{What to Measure?}
		\begin{itemize}
			\item Who has the best average reaction time? \textit{You can use R's \texttt{mean} function.}\pause
			\item Is there any problem with using the above measure?\pause
			\item To begin with, should we even measure times themselves?\pause
			\begin{itemize}
				\item Maybe measure absolute distance from the target value? Thus, opposite errors do not cancel out.\pause
			\end{itemize}
			\item Next, how reliable is mean value as an indicative measure for a particular sample?
		\end{itemize}
	\end{frame}

	\begin{frame}{Okay, any other ideas?}
		\begin{itemize}
			\item A mean value cannot always represent a sample since, both samples below share the same mean value:
			\begin{align*}
				&5, 5, 5, 5, 5, 5, 5, 5, 5, 5, 5, 5, 5, 5, 5, 5, 5, 5, 5, 5, 5\\
				&-2026, -1000, -500, -200, 5, 200, 500, 1000, 2026
			\end{align*}
			\item So, one more piece of information might be a sample's range, $R$:
			\[R=\max x_i-\min x_i.\]
		\end{itemize}
	\end{frame}

	\begin{frame}{More on Mean Values}
		\begin{itemize}
			\item Mean values are, in general, sensitive to extreme observations, i.e., you can manipulate a mean value by adding an extreme enough value.\pause
			\item This is evident from its very definition:
			\[\overline{x}=\frac{1}{n}\sum_{i=1}^nx_i=\frac{x_1+\cdots+x_n}{n}.\]\pause
			\item Changing, e.g., $x_n$ with $3000x_n$ magnifies the effect of $x_n$ on $\overline{x}$.
		\end{itemize}
	\end{frame}

	\begin{frame}{The Median}
		\begin{itemize}
			\item We can think of a number, $\delta$, which splits the sample into two equal parts:
			\begin{itemize}
				\item Part A contains those observations that are at most equal to $\delta$.
				\item Part B contains those observation that are larger than $\delta$.
			\end{itemize}
			\item So, essentially, sort all observations and look exactly at the middle of this sorted sample.
			\item If the number of observations is odd, then there is exactly one such observation.
			\item If the number of observations is even, then there are two "middle" observations, so we compute their mean value.
		\end{itemize}
	\end{frame}

	\begin{frame}{The Median: An Example}
		\begin{itemize}
			\item Consider the following (sorted) sample:
			\[-2, 0, 3, \textcolor{streamlinegreen}{3}, 5, 7, 2027.\]\pause
			\item Consider now the following sample:
			\[-2027, 10, 20, \textcolor{streamlineorange}{\underbrace{30, 40}_{\frac{30+40}{2}=35}}, 70, 70, 80.\]\pause
			\item Observe how change the colourless observations does not affect the sample median at all!
		\end{itemize}
	\end{frame}
	
	\begin{frame}{A More Fine Grained Measure of Dispersion}
		\begin{itemize}
			\item A good idea to measure the dispersion of a sample would be to get the average value of the distances from the mean.\pause
			\item So, something like that?
			\[\]
		\end{itemize}
	\end{frame}
	
\end{document}